% TOP_PROBLEM: {"id":30007593,"problem_number":"LOCAL-30007593","title":"Sources of asset-price volatility","category_id":21,"category":{"id":21,"name":"economics","display_name":"Economics","description":"Open economic theory statements, published research agendas, and proposed scoped specifications; question provenance and historical priority are qualified per record.","slug":"economics","order_index":21,"created_at":"2026-10-08T00:00:00Z"},"status":"research_question","external_url":null,"legacy_ids":[],"published":false,"statement":"In the explicitly specified setting: How much asset-price volatility comes from rational profit news, discount-rate changes and errors in investors’ beliefs? The mathematical statement above fixes the target, assumptions and scope of this catalogue formulation.","economics":{"original_id":82,"field":"Asset pricing and financial markets","kind":"identification","formalization_basis":"adapted_research_question","source_status":"research_agenda","priority_status":"earliest_found","priority_note":"A 1980 working-paper version and independent LeRoy–Porter work precede or overlap the verified 1981 article. No original-priority claim for the broad decomposition.","earliest_verified_reference":{"authors":["Robert J. Shiller"],"title":"Do Stock Prices Move Too Much to Be Justified by Subsequent Changes in Dividends?","year":1981,"url":"https://www.aeaweb.org/aer/top20/71.3.421-436.pdf","doi":null,"locator":"Opening, p.421; p.423; Section VI, pp.433–434","evidence_summary":"Poses dividend-news variance question under a restricted present-value model."},"current_reference":{"authors":["Pedro Bordalo","Nicola Gennaioli","Rafael La Porta","Matthew OBrien","Andrei Shleifer"],"title":"Long-Term Expectations and Aggregate Fluctuations","year":2024,"url":"https://shleifer.scholars.harvard.edu/sites/g/files/omnuum10626/files/shleifer/files/long-term_expectations_and_aggregate_fluctuations.pdf","doi":"10.1086/729198","locator":"Section VI, pp.341–342 (PDF pp.32–33)","evidence_summary":"Calls for expectation formation, propagation and perceived-risk research."},"formalization_note":"The constant-discount variance bound is established. The new task is a decomposition across richer models; recent rational-dividend-news accounts and behavioral accounts give competing partial answers."},"statusQualification":"Published question or research agenda verified at the cited locator. The mathematical specification is adapted for this catalogue and includes additional assumptions; source publication does not establish first-ever priority or certify this exact model remains unsolved. The constant-discount variance bound is established. The new task is a decomposition across richer models; recent rational-dividend-news accounts and behavioral accounts give competing partial answers.","statusReviewedAt":"2026-10-08"}
% FIRST_POSTED: 2026-10-08
% ENTRY_KIND: statement_only
% !TeX program = lualatex
\documentclass[11pt]{article}
\usepackage{fontspec}
\usepackage[a4paper,margin=25mm]{geometry}
\usepackage{amsmath,amssymb,mathtools}
\usepackage{xurl}
\usepackage[hidelinks]{hyperref}
\newcommand{\sref}[2]{\hyperlink{source:#1:#2}{[S#2]}}
\setlength{\emergencystretch}{3em}
\begin{document}
% problemId:30007593
\section*{Sources of asset-price volatility}\label{30007593}
\subsection{Definitions and mathematical statement}
Fix an asset paying nonnegative dividends \(D_t\), its ex-dividend price \(P_t\), an objective probability law \(\mathbb P\), and the information filtration \((\mathcal F_t)\). Assume the discounted dividend sum is square-integrable. Under the restricted constant-discount, rational-expectations benchmark, let
\[P_t^*=\sum_{k\geq1}\delta^kD_{t+k},\qquad P_t=E_{\mathbb P}[P_t^*\mid\mathcal F_t],\quad0<\delta<1.\]
For stationary, consistently detrended objects this implies \(\operatorname{Var}(P_t)\leq\operatorname{Var}(P_t^*)\); it does not apply automatically to changing discount rates. Enlarge the model to allow stochastic marginal utilities, heterogeneous subjective probability laws and dividend-growth news. The task is to identify the contribution of each mechanism to observed price innovations using dividend realizations, investor-expectation data and independently measured financing conditions. Define the contribution as a counterfactual innovation variance when one mechanism is fixed at its conditional baseline and all remaining equilibrium conditions are solved again. Specify the baseline, detrending and terminal-value assumptions explicitly. Determine whether the joint observations identify these counterfactual shares, or only an identified set. Correlations between components must be reported rather than silently assigned to one channel. A rejection of the restricted benchmark is not a proof of irrationality or of an unsolved variance inequality.

\par\textbf{Formulation and scope.} The constant-discount variance bound is established. The new task is a decomposition across richer models; recent rational-dividend-news accounts and behavioral accounts give competing partial answers.

\par\textbf{Open-question provenance.} An explicit question or research agenda was verified at the source locator. This does not by itself certify that every later version remains unresolved \sref{30007593}{1}.

\par\textbf{Historical priority.} A 1980 working-paper version and independent LeRoy–Porter work precede or overlap the verified 1981 article. No original-priority claim for the broad decomposition.

\subsection{Short English statement}
In the explicitly specified setting: How much asset-price volatility comes from rational profit news, discount-rate changes and errors in investors’ beliefs? The mathematical statement above fixes the target, assumptions and scope of this catalogue formulation.

\subsection{Sources}
\begin{itemize}\raggedright
\item[\textnormal{[S1]}]\hypertarget{source:30007593:1}{} Robert J. Shiller. 1981. Do Stock Prices Move Too Much to Be Justified by Subsequent Changes in Dividends?. Opening, p.421; p.423; Section VI, pp.433–434.
\url{https://www.aeaweb.org/aer/top20/71.3.421-436.pdf}.
{\small\textit{Earliest verified question/agenda reference; historical priority qualified below. Poses dividend-news variance question under a restricted present-value model.}}

\item[\textnormal{[S2]}]\hypertarget{source:30007593:2}{} Pedro Bordalo, Nicola Gennaioli, Rafael La Porta, Matthew OBrien, Andrei Shleifer. 2024. Long-Term Expectations and Aggregate Fluctuations. Section VI, pp.341–342 (PDF pp.32–33).
\url{https://shleifer.scholars.harvard.edu/sites/g/files/omnuum10626/files/shleifer/files/long-term_expectations_and_aggregate_fluctuations.pdf}.
DOI: 10.1086/729198.
{\small\textit{Supporting or current literature; see evidence qualification. Calls for expectation formation, propagation and perceived-risk research.}}
\end{itemize}
\end{document}
